\subsection{RQ3: Resources and prerequisites for assessment}
\label{sec:rq3}
\subsubsection{Real reference data and annotations}
\draftnote{Map metrics to unlabelled real images, task labels, subgroup
metadata, paired images or generator training data. Separate generation,
metric development, validation and testing data. Discuss representativeness,
sample size, independent subjects/scenes and uncertainty. State what
cannot be concluded without real data.}

\subsubsection{Statistical priors and domain expertise}
\draftnote{Identify required knowledge of class prevalence, operational
conditions, sensors, physical constraints, rare events and group definitions.
Record each prior's source and uncertainty. Distinguish deployment
prevalence from a balanced evaluation set. Identify expert judgement needs.}

\subsubsection{Trained networks and computation}
\draftnote{Separate pretrained feature extractors for distributional
metrics from task networks trained for utility assessment. Record
checkpoints, training domains, preprocessing, suitability and compute.
Identify network-free measures and domains where generic features may fail.}

\subsubsection{Reproducibility and assessment protocol}
\draftnote{Map each chosen metric to minimum inputs, software/version,
sample-size considerations, budget and limitations. Include generator
settings, seeds, dataset versions, baselines, uncertainty and human-rating
protocols where relevant. Answer RQ3, including resource-limited cases.}
